\chapter{Maximum likelihood and expectation--maximisation}
\label{ch:maximum-likelihood-em}

The paper's learning method is expectation--maximisation (EM) for the marginal
likelihood of noised sequences, with the exact E-step of \cref{ch:factor-graphs-bp}.
This chapter supplies the statistical background in the order a reader needs it: what a
likelihood and a maximum-likelihood estimator are and what consistency and efficiency
mean; what changes when the model has latent variables; Jensen's inequality; the EM
algorithm derived twice, as coordinate ascent on a free-energy bound (Neal and Hinton)
and in the original form of Dempster, Laird and Rubin, with the monotonicity proof;
its generalisations (generalised EM, expectation--conditional maximisation) and why
they keep the monotone ascent; its convergence rate and the fraction of missing
information; stopping rules and the unit of the paper's tolerance; and why EM slows
when the data carry little information about a coordinate, which is the mechanism the
paper's Section 4 invokes when it says the iteration budget binds at high noise.

\section{Likelihood, maximum likelihood, consistency, efficiency}

\subsection{Definitions}

A \term{statistical model} is a family of densities $\{p_\params:\params\in\Theta\}$ on a
sample space, indexed by a parameter in an open set $\Theta\subseteq\Reals^P$. Given
observations $x^1,\dots,x^\nseq$ drawn independently from $p_{\params_0}$ for some
unknown \term{true parameter} $\params_0$, the \term{likelihood} is the joint density as
a function of the parameter, and its logarithm is the \term{log-likelihood}
\begin{equation}
  \Lik(\params)=\sum_{\mu=1}^{\nseq}\log p_\params(x^\mu).
  \label{eq:ml-loglik}
\end{equation}
The \term{maximum-likelihood estimator} (MLE) is any maximiser
$\hat\params=\argmax_\params\Lik(\params)$. Its gradient $\nabla_\params\log p_\params(x)$
is the \term{score} of the model (not to be confused with the score of a density in the
diffusion sense, which is a gradient in $x$; the paper uses ``score'' for the latter and
``Fisher's identity'' for the former); the covariance of the score under $p_\params$ is
the \term{Fisher information} $I(\params)$, treated in the compendium's chapter on Fisher
information and the Cram\'er--Rao bound.

\subsection{Consistency and efficiency in one page}

Two facts justify maximum likelihood. Both are stated for a \term{regular} model: the
densities have a common support, are twice continuously differentiable in $\params$,
differentiation under the integral is allowed, the information matrix is finite and
non-singular, and the parameter is \term{identifiable}, $p_\params=p_{\params'}$ only if
$\params=\params'$.

\emph{Consistency.} By the law of large numbers, $\Lik(\params)/\nseq\to\E_{\params_0}\log p_\params(x)$
for every $\params$, and the limit equals $\E_{\params_0}\log p_{\params_0}(x)-\KL(p_{\params_0}\|p_\params)$.
The Kullback--Leibler divergence is non-negative and zero only at $p_\params=p_{\params_0}$
(this is Gibbs' inequality, itself a consequence of Jensen's inequality below), so the
limiting objective has a unique maximum at $\params_0$ under identifiability. With a
uniform law of large numbers on compact sets, the maximiser of the sample objective
converges to the maximiser of the limit: $\hat\params\to\params_0$ in probability. This
is Wald's argument.

\emph{Asymptotic normality and efficiency.} Expand the score equation
$\nabla\Lik(\hat\params)=0$ around $\params_0$:
$0=\nabla\Lik(\params_0)+\nabla^2\Lik(\bar\params)(\hat\params-\params_0)$ for some
intermediate $\bar\params$. The first term is a sum of $\nseq$ independent zero-mean
vectors (the score has mean zero: $\E_\params\nabla\log p_\params=\nabla\int p_\params=0$),
so by the central limit theorem $\nabla\Lik(\params_0)/\sqrt{\nseq}\to\normal{0}{I(\params_0)}$;
the second factor divided by $\nseq$ converges to $-I(\params_0)$ (the information
identity $\E[-\nabla^2\log p]=\E[\nabla\log p\,\nabla\log p^{\!\top}]$, from differentiating
$\int\nabla p=0$ once more). Hence
\begin{equation}
  \sqrt{\nseq}\,(\hat\params-\params_0)\ \to\ \normal{0}{I(\params_0)^{-1}}.
  \label{eq:ml-asymp}
\end{equation}
The Cram\'er--Rao bound says no unbiased estimator has covariance below $I(\params_0)^{-1}/\nseq$;
an estimator attaining \cref{eq:ml-asymp} is \term{asymptotically efficient}, and the
MLE is one. The local asymptotic minimax theorem of H\'ajek and Le Cam, which the
paper's sample-complexity corollary invokes, is the sharp form: over shrinking neighbourhoods
$|\params'-\params_0|\le c/\sqrt\nseq$ no estimator sequence does better than
$1/I(\params_0)$ in the limit, and the MLE attains it. Three caveats travel with these
statements and the paper repeats them: they are limits in $\nseq$ at fixed model, they
say nothing about a finite optimisation budget, and for a bounded coordinate the local
statement is empty once $1/\sqrt{\nseq I}$ exceeds the coordinate's range.

\section{Latent-variable models}

A \term{latent-variable model} specifies a joint density $p_\params(x,z)$ of an observed
variable $x$ and an unobserved (\term{latent}, \term{hidden}, \term{missing}) variable
$z$, and the likelihood of the data is the \term{marginal likelihood}
\begin{equation}
  p_\params(x)=\int p_\params(x,z)\dd z.
  \label{eq:ml-marginal}
\end{equation}
The joint is usually simple, often in an exponential family whose log-likelihood is
linear in a few sufficient statistics of $(x,z)$; the marginal is usually not, because
the integral destroys the structure. In the paper $x$ is a noised sequence, $z$ is the
clean chain, the joint is $\Pz(\va)\prod_i\like(x_i\mid a_i)$, and the integral over the
chain is what the two sweeps of \cref{ch:factor-graphs-bp} compute. The
\term{complete-data log-likelihood} $\log p_\params(x,z)$ is the object EM works with:
were $z$ observed, maximising it would be easy; it is not, so EM replaces it by its
conditional expectation given the data.

\section{Jensen's inequality}

\begin{lemma}[Jensen]
\label{lem:ml-jensen}
Let $\phi$ be a concave function on an interval $J$ and $Y$ a random variable with
values in $J$ and finite expectation. Then $\phi(\E Y)\ge\E\phi(Y)$, with equality if
$\phi$ is strictly concave only when $Y$ is almost surely constant.
\end{lemma}
\begin{proof}
A concave function lies below every one of its supporting lines: at the point
$m=\E Y$ there is a slope $s$ with $\phi(y)\le\phi(m)+s(y-m)$ for all $y\in J$ (for a
differentiable $\phi$, $s=\phi'(m)$; in general any element of the superdifferential).
Take expectations: $\E\phi(Y)\le\phi(m)+s(\E Y-m)=\phi(m)$. For strict concavity the
inequality $\phi(y)<\phi(m)+s(y-m)$ is strict for $y\ne m$, so equality of expectations
forces $Y=m$ almost surely.
\end{proof}

With $\phi=\log$, which is strictly concave on $(0,\infty)$, two consequences follow.
\term{Gibbs' inequality}: for densities $p,q$ with $q\ll p$,
$\int p\log(q/p)\le\log\int p\,(q/p)=\log\int q\le0$, so $\KL(p\|q)=\int p\log(p/q)\ge0$
with equality only at $p=q$. And the \term{evidence lower bound} of the next section.

\section{Expectation--maximisation}

\subsection{The free-energy view}

Fix the data $x$ and consider an arbitrary density $q(z)$ over the latent variable.
Write the log marginal likelihood as an expectation under $q$ and apply
\cref{lem:ml-jensen}:
\begin{equation}
  \log p_\params(x)=\log\int q(z)\frac{p_\params(x,z)}{q(z)}\dd z
  \ \ge\ \int q(z)\log\frac{p_\params(x,z)}{q(z)}\dd z
  =\underbrace{\E_q\log p_\params(x,z)}_{\text{energy}}+\underbrace{H(q)}_{\text{entropy}}
  =:F(q,\params).
  \label{eq:ml-elbo}
\end{equation}
The bound $F$ is the \term{free energy} (negative, in the physicists' sign) or
\term{evidence lower bound}; it is exact in the sense that
\begin{equation}
  \log p_\params(x)-F(q,\params)=\KL\bigl(q\,\|\,p_\params(\cdot\mid x)\bigr)\ \ge0,
  \label{eq:ml-gap}
\end{equation}
which one checks by writing $p_\params(x,z)=p_\params(x)p_\params(z\mid x)$ inside the
integral. So $F(q,\params)\le\log p_\params(x)$ for every $q$, with equality exactly when
$q$ is the posterior $p_\params(\cdot\mid x)$.

The EM algorithm of Neal and Hinton (1998) is \term{coordinate ascent} on $F$:
\begin{itemize}
\item \term{E-step}: with $\params^{(k)}$ fixed, maximise $F$ over $q$. By
\cref{eq:ml-gap} the maximiser is the posterior $q^{(k+1)}=p_{\params^{(k)}}(\cdot\mid x)$,
and at it $F=\log p_{\params^{(k)}}(x)$.
\item \term{M-step}: with $q^{(k+1)}$ fixed, maximise $F$ over $\params$. The entropy
does not depend on $\params$, so this is maximising the expected complete-data
log-likelihood $\E_{q^{(k+1)}}\log p_\params(x,z)$; call the maximiser $\params^{(k+1)}$.
\end{itemize}
Each step does not decrease $F$, and after every E-step $F$ equals the log-likelihood.
Hence $\log p_{\params^{(k+1)}}(x)=F(q^{(k+2)},\params^{(k+1)})\ge F(q^{(k+1)},\params^{(k+1)})\ge F(q^{(k+1)},\params^{(k)})=\log p_{\params^{(k)}}(x)$:
the marginal likelihood never decreases. The view also explains partial versions of
either step: any $q$ that increases $F$ and any $\params$ that increases $F$ preserve
the monotonicity, because $F$ is the single objective both steps ascend.

\subsection{The Dempster--Laird--Rubin view and the Q-function}

Define the \term{Q-function}
\begin{equation}
  \Qfun(\params\mid\params')=\E_{p_{\params'}(z\mid x)}\bigl[\log p_\params(x,z)\bigr]
  =\int p_{\params'}(z\mid x)\log p_\params(x,z)\dd z,
  \label{eq:ml-Q}
\end{equation}
the expected complete-data log-likelihood at $\params$ under the posterior computed at
$\params'$. The E-step computes whatever expectations $\Qfun(\cdot\mid\params^{(k)})$
needs; the M-step sets $\params^{(k+1)}=\argmax_\params\Qfun(\params\mid\params^{(k)})$. The
original monotonicity proof (Dempster, Laird and Rubin, 1977) is the following.

\begin{theorem}[Monotone ascent of EM]
\label{thm:ml-monotone}
If $\Qfun(\params^{(k+1)}\mid\params^{(k)})\ge\Qfun(\params^{(k)}\mid\params^{(k)})$, then
$\log p_{\params^{(k+1)}}(x)\ge\log p_{\params^{(k)}}(x)$, with equality only if the
posteriors at the two parameters coincide and the $\Qfun$ values are equal.
\end{theorem}
\begin{proof}
Write $\log p_\params(x)=\log p_\params(x,z)-\log p_\params(z\mid x)$ for any $z$ and take
the expectation under $p_{\params'}(z\mid x)$, which leaves the left side unchanged:
\begin{equation}
  \log p_\params(x)=\Qfun(\params\mid\params')-\mathcal H(\params\mid\params'),\qquad
  \mathcal H(\params\mid\params')=\E_{p_{\params'}(z\mid x)}\log p_\params(z\mid x).
  \label{eq:ml-decomp}
\end{equation}
By Gibbs' inequality, $\mathcal H(\params\mid\params')\le\mathcal H(\params'\mid\params')$ for every
$\params$ (the cross-entropy of the posterior at $\params'$ against the posterior at
$\params$ is largest when the two agree). Subtracting \cref{eq:ml-decomp} at
$\params=\params^{(k+1)}$ from itself at $\params=\params^{(k)}$, both with
$\params'=\params^{(k)}$,
\[
  \log p_{\params^{(k+1)}}(x)-\log p_{\params^{(k)}}(x)
  =\bigl[\Qfun(\params^{(k+1)}\mid\params^{(k)})-\Qfun(\params^{(k)}\mid\params^{(k)})\bigr]
  +\bigl[\mathcal H(\params^{(k)}\mid\params^{(k)})-\mathcal H(\params^{(k+1)}\mid\params^{(k)})\bigr],
\]
and both brackets are non-negative.
\end{proof}

The theorem asks only that $\Qfun$ \emph{increase}, not that it be maximised. This is the
definition of \term{generalised EM} (GEM): any M-step that ascends $\Qfun$ preserves
the monotone increase of the likelihood. A useful special case is
\term{expectation--conditional maximisation} (ECM, Meng and Rubin, 1993): partition the
parameter into blocks $\params=(\params_1,\dots,\params_B)$ and replace the M-step by a
sweep of \term{conditional maximisations}, each maximising $\Qfun$ over one block with
the others held at their most recent values. Each conditional maximisation cannot
decrease $\Qfun$ (it maximises it over a subset containing the current point), so the
sweep ascends $\Qfun$ and \cref{thm:ml-monotone} applies; several sweeps, or a
different order, are equally admissible. Meng and Rubin also show that if every
conditional maximiser is unique and the blocks together span the parameter space in a
suitable sense, the fixed points of ECM are the stationary points of the likelihood,
as for EM. The paper's mixture-innovation M-step is an ECM sweep, block by block,
which is why the paper calls it ``a generalised EM whose ascent is monotone'' and
cites Meng and Rubin.

\subsection{Fixed points are stationary points}

At a fixed point $\params^\ast$ of EM, $\params^\ast$ maximises $\Qfun(\cdot\mid\params^\ast)$,
so $\nabla_\params\Qfun(\params\mid\params^\ast)|_{\params=\params^\ast}=0$. Differentiating
\cref{eq:ml-decomp} in $\params$ at $\params=\params'$ shows
$\nabla_\params\mathcal H(\params\mid\params')|_{\params=\params'}=\E_{p_{\params'}(z\mid x)}\nabla\log p_{\params'}(z\mid x)=0$
(the posterior score has mean zero under the posterior), hence
\begin{equation}
  \nabla_\params\log p_\params(x)=\nabla_\params\Qfun(\params\mid\params')\big|_{\params'=\params}
  =\E_{p_\params(z\mid x)}\bigl[\nabla_\params\log p_\params(x,z)\bigr].
  \label{eq:ml-fisher}
\end{equation}
This is \term{Fisher's identity}: the gradient of the marginal log-likelihood is the
posterior expectation of the complete-data gradient. It says that fixed points of EM
are stationary points of the likelihood, and, more usefully for the paper, that the
E-step delivers the exact gradient of the marginal likelihood as a by-product.
\Cref{ch:em-through-noise} derives it again directly, in the form the paper uses.

\section{Convergence rate and the fraction of missing information}

EM defines a map $\params^{(k+1)}=\mathcal M(\params^{(k)})$. Near a fixed point $\params^\ast$
with a non-singular Hessian, $\params^{(k+1)}-\params^\ast\approx J(\params^{(k)}-\params^\ast)$
with $J=\nabla\mathcal M(\params^\ast)$ the Jacobian of the map, so the error contracts
linearly with rate equal to the largest eigenvalue of $J$. Dempster, Laird and Rubin
identified $J$. Define the \term{complete information}
$I_{\mathrm{com}}=-\nabla^2_\params\Qfun(\params\mid\params')|_{\params=\params'=\params^\ast}$,
the expected negative Hessian of the complete-data log-likelihood under the posterior,
and the \term{observed information} $I_{\mathrm{obs}}=-\nabla^2\log p_{\params^\ast}(x)$. Then
\begin{equation}
  J=I_{\mathrm{mis}}\,I_{\mathrm{com}}^{-1},\qquad I_{\mathrm{mis}}=I_{\mathrm{com}}-I_{\mathrm{obs}},
  \qquad\text{so}\quad J=\Id-I_{\mathrm{obs}}I_{\mathrm{com}}^{-1}.
  \label{eq:ml-rate}
\end{equation}
The matrix $I_{\mathrm{mis}}$ is the \term{missing information}, the expected information
that the latent variable would have added, and $J$ is the \term{fraction of missing
information}. The derivation is one line from the M-step's stationarity condition:
$\nabla_\params\Qfun(\mathcal M(\params')\mid\params')=0$ for all $\params'$; differentiate
in $\params'$ at the fixed point to get
$-I_{\mathrm{com}}J+\nabla_{\params'}\nabla_\params\Qfun=0$, and the mixed derivative equals
$I_{\mathrm{com}}-I_{\mathrm{obs}}$ by differentiating \cref{eq:ml-decomp} twice (the
second derivative of $\mathcal H$ in its first argument at the fixed point is minus the
posterior information, and the identity rearranges to \cref{eq:ml-rate}).

The reading is the important part. When the data determine the parameter almost as
well as the complete data would, $I_{\mathrm{obs}}\approx I_{\mathrm{com}}$, $J\approx0$ and
EM converges in a few iterations. When the observed information is a small fraction of
the complete information, the eigenvalues of $J$ approach one and EM crawls: each
iteration removes only the fraction $1-\lambda_{\max}(J)$ of the remaining error, and
the number of iterations to reduce the error by a factor $e$ is about
$1/(1-\lambda_{\max}(J))$. In a multi-parameter problem the rate is direction-dependent:
a coordinate about which the data are nearly uninformative has its own eigenvalue
near one and converges last, however fast the others settle. This is the observation
the next section, and the paper, build on.

\section{Stopping rules}

An EM run is stopped by a \term{tolerance} on the increase of the log-likelihood, by an
\term{iteration cap}, or by an external \term{selection rule}. The paper uses all three
and the compendium should say what each measures.

\paragraph{Units.} The total log-likelihood \cref{eq:ml-loglik} of $\nseq$ sequences of
$\nsites$ sites is a sum over $\nseq(\nsites-1)$ transitions (plus the first sites), so
it grows with the dataset, and so does its increment per iteration. A tolerance in
absolute nats means something different at every sample size. A relative tolerance
$|\Delta\Lik|\le\tau|\Lik|$ is scale-free, because the increment and the total are sums
over the same transitions; the docstring of \code{fit\_em} in
\code{src/markov\_score/em.py} records that this was measured to be invariant to the
sample size to every digit, so the relative rule was not wrong. What the paper's code
does instead is to express the increment in \term{nats per transition},
$\Delta\Lik/\bigl(\nseq(\nsites-1)\bigr)$, and to stop when it lies in $[0,\tau]$ with
$\tau=10^{-9}$. The unit buys interpretability: a constant threshold in nats per
transition means the same thing at every sample size and every noise level, whereas a
relative threshold tightens as $|\Lik|$ falls along the run and changes its meaning
with the level (at high noise $|\Lik|$ per transition is close to the entropy of a
standard Gaussian and the increments are tiny in comparison).

\paragraph{Sign.} An exact EM cannot decrease the likelihood. A stopping test on
$|\Delta\Lik|$ accepts a small \emph{decrease} as convergence and thereby hides a defect
(a wrong M-step, or a generalised step that overshot). The code's test is signed:
convergence requires $0\le\Delta\Lik\le\tau$ per transition; any decrease is recorded as
a \code{monotone\_violation} and is a bug report, not a stop. A run that exhausts its
iteration cap without meeting the signed test is labelled \term{censored}, and the
label travels with the fitted kernel so that a comparison of estimators can tell
whether it compares converged runs. The returned kernel is always the iterate whose
likelihood is the last trace entry: the loop scores each proposal after the M-step that
produced it. An earlier version returned a kernel one update beyond anything evaluated,
which made the reported likelihood belong to the previous iterate; the tests in
\code{tests/test\_em\_bp.py} and \code{tests/test\_gradient\_trace\_alignment.py} pin
this alignment.

\paragraph{Selection.} At the budgets of the paper's experiments, the per-transition
increments EM takes are far above $10^{-9}$ for hundreds of iterations, so the
tolerance rarely fires and what ends a run is the cap of the \term{checkpoint ladder}
(\cref{ch:em-through-noise}). Which rung is reported is then decided by a validation
bundle, not by the likelihood: the rung with the smallest score error on held-out
sequences. This is \term{early stopping}, and it is a form of regularisation, because
successive EM iterates move from the initialisation towards the maximiser and an early
rung is a shrunken estimate. The paper's Section 4 reads its own fits this way: at
high noise the validation-selected iterate is a shrinkage estimator whose shape
collapses towards a Gaussian innovation. The paper's Appendix E (E20) checks that
selecting the rung by held-out log-evidence instead of by the validation score error
picks the same rung in most seeds, so the selection needs no oracle.

\section{Why EM slows when information is scarce}
\label{sec:ml-slow}

Combine \cref{eq:ml-rate} with the paper's order-by-order decay theorem. In the
noised-sequence problem
the complete data are the clean chains, and the complete information about the kernel
parameters is of order $\nseq(\nsites-1)$ times a constant that does not depend on the
noise level, since it is the information the clean transitions carry. The observed
information is $\nseq$ times the per-sequence Fisher information $I_t(\params)$, which
decays as $e^{-4t}$ for a covariance coordinate and $e^{-8t}$ for a shape coordinate.
Hence the fraction of missing information along each coordinate,
\begin{equation}
  \lambda(\params_j)\approx1-\frac{I_t(\params_j)}{I_{\mathrm{com}}(\params_j)},
  \label{eq:ml-fraction}
\end{equation}
approaches one as $t$ grows, and does so faster for the shape coordinate. Two
consequences follow, both visible in the paper. First, the number of iterations to
converge grows like $1/(1-\lambda)$, so at high noise the ladder's cap is reached
before the tolerance fires: the paper reports that at the two highest levels of its
learning experiment the validation-selected rung is the last of the ladder in most
cells, ``whether the excess over the reference at those two levels is information or
budget the run cannot say'', and registers a control with a sixfold longer ladder
(E06c). Second, at any fixed noise level the correlation coordinate (degree two)
converges in far fewer iterations than the shape coordinate (degree four), because
its missing fraction is smaller by roughly $e^{-4t}$ in the information ratio; the
paper's ``convergence asymmetry'' (correlation settles in tens of updates, shape in
hundreds) is \cref{eq:ml-fraction} read coordinate by coordinate.
\Cref{ch:em-through-noise} adds the second mechanism specific to the mixture kernel,
the inner latent label, which slows the shape further.

\section{A worked example: one mixing weight}

Let $y_1,y_2,y_3=-1.5,\ 0.2,\ 1.0$ be three observations from the two-component
mixture $p_\pi(y)=\pi\,\phi(y-1)+(1-\pi)\,\phi(y+1)$ with $\phi$ the standard Gaussian
density and the mixing weight $\pi$ unknown. The latent variable is the component
label $z_i\in\{0,1\}$; the complete-data log-likelihood is
$\sum_i[z_i\log\pi+(1-z_i)\log(1-\pi)]+\text{const}$, linear in the $z_i$, so the E-step
needs only the \term{responsibilities} $r_i=\E[z_i\mid y_i]=\pi\phi(y_i-1)/p_\pi(y_i)$, and
the M-step maximises $\sum_i[r_i\log\pi+(1-r_i)\log(1-\pi)]$, giving
$\pi^{\mathrm{new}}=\tfrac13\sum_ir_i$.

The densities are
\[
  \phi(y_i-1)=(0.01753,\,0.28969,\,0.39894),\qquad
  \phi(y_i+1)=(0.35207,\,0.19419,\,0.05399).
\]
Start at $\pi^{(0)}=0.5$. The log-likelihood is
$\sum_i\log[0.5\phi(y_i-1)+0.5\phi(y_i+1)]=-4.59273$. The responsibilities are
$r=(0.04743,\,0.59869,\,0.88080)$ and the M-step gives $\pi^{(1)}=0.50897$. The
log-likelihood rises to $-4.59199$; the new responsibilities are
$(0.04907,\,0.60728,\,0.88451)$ and $\pi^{(2)}=0.51362$, with log-likelihood
$-4.59180$; then $\pi^{(3)}=0.51602$,
$\pi^{(4)}=0.51725$, $\pi^{(5)}=0.51789$, with log-likelihoods $-4.59174$, $-4.59173$,
$-4.59173$: monotone, as \cref{thm:ml-monotone} requires, and slowing.

The maximiser is $\pi^\ast=0.51855$. The successive errors $|\pi^{(k)}-\pi^\ast|$ shrink
by the factor $0.514$, $0.514$, $0.513$, \dots\ per iteration. \Cref{eq:ml-rate}
predicts this number: at $\pi^\ast$ the complete information is
$I_{\mathrm{com}}=\sum_i[r_i/\pi^2+(1-r_i)/(1-\pi)^2]=12.02$, the observed information is
minus the second derivative of the log-likelihood, $I_{\mathrm{obs}}=5.86$ (a finite
difference suffices), and $1-I_{\mathrm{obs}}/I_{\mathrm{com}}=0.513$. Half the information
about $\pi$ is missing, and EM removes half of the remaining error per iteration. Had
the two components been closer (means $\pm0.2$ instead of $\pm1$), the responsibilities
would be near $\tfrac12$ for every observation, the observed information a small
fraction of the complete one, and the rate near one: the situation of
\cref{sec:ml-slow}.

\where{The paper, Section 4 (``expectation--maximisation around the same exact
recursion'', generalised EM with conditional-maximisation sweeps, the convergence
asymmetry, the iteration budget at high noise), the sample-complexity corollary and its
caveats (local
asymptotic minimax, the MLE attains it), Appendix C (the checkpoint ladder and
validation selection), Appendix E (E20). Code: \code{src/markov\_score/em.py}, whose
\code{fit\_em} carries the docstring on the stopping rule and whose \code{EMTrace}
records \code{monotone\_violation}, \code{converged} and \code{stop\_reason}, with
\code{q\_value} and \code{q\_gradient} beside them; and the ECM sweeps of the mixture
kernel's \code{m\_step} in \code{src/markov\_score/kernels.py}. Tests: monotonicity,
the signed stopping test, censored against converged and the stationarity of the fixed
point in \code{tests/test\_em\_bp.py}; the iterate-to-trace alignment in
\code{tests/test\_gradient\_trace\_alignment.py}.}

\paragraph{References.}
A.~P.~Dempster, N.~M.~Laird and D.~B.~Rubin (1977), ``Maximum likelihood from incomplete
data via the EM algorithm'', \emph{Journal of the Royal Statistical Society B} 39(1).
R.~M.~Neal and G.~E.~Hinton (1998), ``A view of the EM algorithm that justifies
incremental, sparse, and other variants'', in \emph{Learning in Graphical Models}, ed.\
M.~I.~Jordan, Kluwer.
X.-L.~Meng and D.~B.~Rubin (1993), ``Maximum likelihood estimation via the ECM algorithm:
a general framework'', \emph{Biometrika} 80(2).
C.~F.~J.~Wu (1983), ``On the convergence properties of the EM algorithm'', \emph{Annals of
Statistics} 11(1).
A.~W.~van der Vaart (1998), \emph{Asymptotic Statistics}, Cambridge University Press,
Chapters 5, 7 and 8.
O.~Capp\'e, E.~Moulines and T.~Ryd\'en (2005), \emph{Inference in Hidden Markov Models},
Springer, Chapter 10.
